\section{Migration Strategy}
\label{sec:migration}

\begin{diagram}
 PHASE 0: Observe               Legacy vehicle + Alfred Probe
                                 (Section 4; passive only)
      |
      v
 PHASE 1: Understand            Vehicle Knowledge Graph populated:
                                 machine graph, ECU map, signal map
                                 (Section 7)
      |
      v
 PHASE 2: Control selected      Alfred API -> Alfred Gateway ->
          functions             legacy network adapter (Section 10.4)
                                 -- non-safety-critical scope only
                                 (Section 9.3)
      |
      v
 PHASE 3: Replace selected      old body ECU removed -> Alfred E2B
          ECUs                  node installed in its place, using
                                 the SAME connector/harness point
                                 where feasible (Section 3.2)
      |
      v
 PHASE 4: Collapse network      CAN / LIN islands, now serving fewer
          architecture          and fewer legacy ECUs, are progressively
                                 retired in favor of Ethernet edge
                                 architecture (Section 3.1/3.2)
      |
      v
 PHASE 5: Alfred-native         central compute + Ethernet backbone +
          vehicle               E2B nodes (Section 3)
\end{diagram}

\subsection{Why this is a low-risk entry motion}

\begin{itemize}
\item \textbf{Phase 0/1 require zero vehicle modification and zero actuation risk} --- a customer can authorize probe installation on a development/test vehicle with essentially the same risk profile as installing a data logger, which is a far smaller approval/procurement hurdle than "let us reprogram your ECUs."
\item \textbf{Phase 2 stays within the non-safety-critical allowlist} (Section~\ref{sec:active-control}) by design, so the first revenue-relevant capability Alfred demonstrates (e.g., programmatic control of seats/mirrors/HVAC for a fleet customer's ergonomics or comfort feature) never requires the customer to accept safety-case risk on Alfred's word alone.
\item \textbf{Phase 3 targets are chosen by the customer, one ECU at a time}, using the Vehicle Knowledge Graph's confidence scores and the ECU's programmability assessment (Section~\ref{sec:deployment-reprogramming}) to pick the lowest-risk, highest-value replacement first --- typically a body-function ECU with a well-characterized command surface and no safety classification.
\item \textbf{Phase 4/5 are economically self-reinforcing once started}: each ECU replaced with an E2B node removes a CAN/LIN drop, which reduces bus load and complexity for the \emph{remaining} legacy ECUs' traffic, which in turn makes each subsequent migration incrementally easier to validate (less contention, cleaner captures) --- the architecture gets easier to finish the more of it is finished, which is an unusual and favorable property for a phased hardware migration.
\end{itemize}

\subsection{Entry-then-expand as the commercial mechanism}

This phased structure is what lets Alfred enter a customer account through a single low-risk subsystem (a lighting or seat-control pilot, say) and expand its footprint over the following 12--36 months without ever requiring the customer to make a single large architecture bet up front. Each phase transition is a new, independently justifiable purchase decision (a Probe engagement, then a Gateway deployment, then a handful of E2B nodes, then a broader retrofit program) rather than one all-or-nothing platform commitment --- which matters commercially because automotive/industrial buyers are structurally averse to the latter and structurally comfortable with the former.
